% ==============================================================================
% ================================ Zusammenfassung ==============================
% ==============================================================================

\chapter{Zusammenfassung und Reflexion}
\label{chap:zusammenfassung}
\thispagestyle{fancy}

Zusammenfassend lässt sich festhalten, dass die Transformation der \acs{SMT} zur \textit{Smart Factory} auf einer ganzheitlichen Vernetzung und der intelligenten Nutzung von Echtzeitdaten basiert. Hier sind die Kernpunkte des Beratungskonzepts:

\begin{itemize}[leftmargin=*]
    \item \textbf{Cyber-Physische Systeme \textit{(\acs{CPS})}}: Gemäß Jeanetta Zotter und Andreas Metzler \cite[4-5]{2} bilden \acs{CPS} über fünf Ebenen das technologische Rückgrat, um physische Abläufe digital abzubilden und autonom zu steuern.
    
    \item \textbf{Software-Integration}: Wie Marion Steven und Jan Niklas Dörseln \cite[11]{1} ausführen, ist die durchgängige Integration aller Ebenen der \textit{Automatisierungspyramide}, von der \acs{ERP}- bis zur \acs{MES}-Ebene, essenziell für die nötige Transparenz.
    
    \item \textbf{Vorausschauende Wartung}: Laut Jeanetta Zotter und Andreas Metzler \cite[39]{2} ermöglichen Echtzeit-Sensordaten den entscheidenden Übergang von der reaktiven zur vorausschauenden Wartung \textit{(Predictive Maintenance)}, was Stillstände minimiert.
    
    \item \textbf{Cloud \& Big Data}: Marion Steven und Jan Niklas Dörseln \cite[75-76]{1} heben hervor, dass Cloud-Computing die notwendige Skalierbarkeit und Kosteneffizienz bietet, um große Datenmengen \textit{(Big Data)} ohne hohe Hardware-Investitionen zu verarbeiten.
    
    \item \textbf{Sichere Vernetzung}: Zur Gewährleistung der Stabilität empfiehlt sich ein hochverfügbares Netzwerkdesign nach dem \textit{\acs{IEC} 62443}-Standard, der, so Zotter und Metzler, den maßgeblichen Leitfaden für die industrielle Sicherheit darstellt \cite[73]{2}.
    
    \item \textbf{Nachhaltigkeit}: Wie Zotter und Metzler \cite[79]{2} dargelegt haben, führt das \glqq Green \acs{IIoT}\grqq{} \acs{SMT} durch ressourceneffiziente Prozesse und intelligentes Energiemanagement konsequent zur nachhaltigen Produktion.
\end{itemize}

Abschließend bleibt zu reflektieren, dass der entscheidende Erfolgsfaktor für \acs{SMT} nicht allein in der Technologie liegt, sondern im Wandel von der reaktiven Problembehebung zur proaktiven Fehlervermeidung \textit{(\acs{vgl.} Steven und Dörseln \cite[11]{1})}. Dies sichert langfristig die Qualität und Wettbewerbsfähigkeit bei Losgröße 1.

% ==============================================================================
% ================================ Zusammenfassung ==============================
% ==============================================================================